\documentclass[a4paper,12pt]{article}
\usepackage{iftex}
\ifPDFTeX
  \usepackage{cmap}
  \usepackage[T2A]{fontenc}
  \usepackage[utf8]{inputenc}
  \usepackage[russian]{babel}
  \usepackage{upquote}
\else
  \usepackage{fontspec}
  \usepackage{polyglossia}
  \setdefaultlanguage{russian}
  \setmainfont{Times New Roman}
  \setsansfont{Arial}
  \setmonofont{Consolas}
  \usepackage{upquote}
\fi
\usepackage{indentfirst}
\usepackage{listings}
\usepackage{listingsutf8}
\usepackage{xcolor}
\usepackage{geometry}
\usepackage{hyperref}
\usepackage{enumitem}
\usepackage{graphicx}
\usepackage{float}
\usepackage{amsmath}
\geometry{top=2cm, bottom=2cm, left=3cm, right=2cm}

\definecolor{codegreen}{rgb}{0,0.6,0}
\definecolor{codegray}{rgb}{0.5,0.5,0.5}
\definecolor{codepurple}{rgb}{0.58,0,0.82}
\definecolor{backcolour}{rgb}{0.97,0.97,0.97}
\lstset{
  basicstyle=\ttfamily\small,
  backgroundcolor=\color{backcolour},
  commentstyle=\color{codegreen},
  keywordstyle=\color{blue},
  numberstyle=\tiny\color{codegray},
  stringstyle=\color{codepurple},
  breaklines=true,
  captionpos=b,
  keepspaces=true,
  numbers=left,
  numbersep=5pt,
  showspaces=false,
  showstringspaces=false,
  showtabs=false,
  tabsize=4,
  columns=fullflexible,
  upquote=true,
  inputencoding=utf8,
  extendedchars=false,
  language=[5.1]Lua,
  morekeywords={Ledbar,new,set,sleep,ap,push,Ev,Timer,callLater,start,stop,unpack,table}
}
\sloppy
\emergencystretch=1em
\setlist[enumerate]{label=\arabic*.,leftmargin=*,itemsep=2pt}

\title{\textbf{Урок №1: Программирование светодиодной индикации и событийная модель}}
\author{Образовательный курс «Геоскан Пионер»}
\date{}
\begin{document}
\maketitle

\section{Введение}
Этот урок вводит в программирование индикации на квадрокоптере «Геоскан Пионер» на языке Lua. Мы начнём с определений, затем — теория, примеры и практические задания.

\section{Толковый словарь}
Нажмите на термин в тексте, чтобы перейти к его определению.
\begin{description}
  \item[\hypertarget{term:module}{Модуль}] часть программы, отвечающая за конкретную задачу и предоставляющая функции для работы. Модуль \hyperlink{term:ledbar}{Ledbar} управляет светодиодами.
  \item[\hypertarget{term:driver}{Драйвер}] программный компонент, который общается с аппаратурой и переводит команды в сигналы для устройств.
  \item[\hypertarget{term:buffer}{Буфер}] область памяти, куда мы записываем значения (например, цвета); драйвер считывает буфер и выводит свечение.
  \item[\hypertarget{term:variable}{Переменная}] именованная «коробочка» со значением. Пример: \verb|x = 10|.
  \item[\hypertarget{term:scope}{Локальная/Глобальная}] локальная переменная объявляется через \verb|local| и видна только внутри текущего блока; глобальная (без \verb|local|) видна во всём скрипте.
  \item[\hypertarget{term:timer}{Таймер}] «будильник», который вызывает указанную функцию периодически или один раз с задержкой.
  \item[\hypertarget{term:period}{Период}] время между двумя вызовами таймера.
  \item[\hypertarget{term:tick}{Тик}] один вызов функции таймером.
  \item[\hypertarget{term:state}{Состояние}] текущее положение логики (например, «красный», «жёлтый», «зелёный»).
  \item[\hypertarget{term:automaton}{Конечный автомат}] модель, где программа находится в одном из состояний и переходит в другое при событиях.
  \item[\hypertarget{term:rgb}{RGB-компоненты}] три числа \verb|R|, \verb|G|, \verb|B|, задающие яркости каналов; каждое в \verb|[0.0, 1.0]|.
  \item[\hypertarget{term:index}{Индекс светодиода}] порядковый номер «лампочки» в линейке; используем \verb|0..24|.
  \item[\hypertarget{term:api}{API}] набор функций модуля (\verb|Ledbar.set|, \verb|Timer.new| и др.).
  \item[\hypertarget{term:ledbar}{Ledbar}] модуль управления RGB-светодиодами квадрокоптера.
\end{description}

\section{Аппаратная часть и модуль Ledbar}
\subsection{Физическая реализация}
На плате установлено \textbf{4 RGB-светодиода}. Каждый имеет три кристалла: красный (R), зелёный (G), синий (B). Комбинация трёх яркостей даёт итоговый цвет.

\subsection{Особенности адресации (Буфер драйвера)}
Важно: \hyperlink{term:ledbar}{Ledbar} работает с \hyperlink{term:buffer}{буфером} \textbf{из 25 элементов}. Поэтому инициализация должна задавать размерность \verb|25|.
\begin{quote}
\textbf{Требование:} \verb|Ledbar.new(25)|. Значение \verb|4| использовать нельзя — диоды не загорятся.
\end{quote}
\begin{lstlisting}[style=lua,caption={Инициализация и установка зелёного цвета}]
local unpack = table.unpack
local ledNumber = 25                  -- размерность буфера
local leds = Ledbar.new(ledNumber)    -- объект управления лентой

local function changeColor(col)
  for i=0, ledNumber-1 do             -- индексы 0..24
    leds:set(i, unpack(col))          -- записываем RGB
  end
end

changeColor({0,1,0})                  -- включаем зелёный
\end{lstlisting}

\section{Теория и интерфейсы API}
\subsection{RGB-модель и диапазоны}
\hyperlink{term:rgb}{RGB} — три управляющих числа: \verb|R| (красный), \verb|G| (зелёный), \verb|B| (синий) в \verb|[0.0, 1.0]|. \verb|0.0| — канал выключен, \verb|1.0| — горит ярко. Смешивая каналы, получаем цвета: \verb|{1,0,0}| (красный), \verb|{0,1,0}| (зелёный), \verb|{0,0,1}| (синий), \verb|{1,1,0}| (жёлтый), \verb|{1,1,1}| (белый), \verb|{0,0,0}| (выключено). Плавная смена значений позволяет создавать анимации.

\subsection{Адресация и буфер}
\hyperlink{term:index}{Индекс} — позиция \verb|0..24|. \hyperlink{term:buffer}{Буфер} — массив из 25 ячеек, где каждая хранит цвет своего индекса. Драйвер считывает весь буфер, поэтому обновляйте все элементы, даже если видны только первые четыре.

\subsection{Модуль Ledbar}
\begin{itemize}
  \item \verb|Ledbar.new(N)| — создаёт объект управления (\hyperlink{term:api}{API}); используйте \verb|25|.
  \item \verb|leds:set(i, r, g, b)| — задаёт цвет индекса \verb|i| (\verb|r,g,b| в \verb|[0.0,1.0]|).
  \item Вспомогательная \verb|changeColor(col)| перезаписывает все индексы.
\end{itemize}
\textbf{Размерность} — число элементов буфера (\verb|25|). \textbf{Индекс} — порядковый номер, начиная с \verb|0|.

\subsection{Таймеры}
\begin{itemize}
  \item \verb|Timer.new(period, fn)| — создаёт \hyperlink{term:timer}{периодический таймер}; \verb|period| — \hyperlink{term:period}{период} в секундах; \verb|fn| вызывается каждый \hyperlink{term:tick}{тик}.
  \item \verb|timer:start()|, \verb|timer:stop()| — запуск/остановка таймера.
  \item \verb|Timer.callLater(delay, fn)| — единичный вызов \verb|fn| через \verb|delay| секунд.
  \item Таймеры и \hyperlink{term:state}{состояния} делайте \hyperlink{term:scope}{глобальными} (без \verb|local|), иначе сборщик мусора их удалит.
\end{itemize}
\textbf{Пример периодического таймера (мигание):}
\begin{lstlisting}[style=lua]
local ledNumber = 25
local leds = Ledbar.new(ledNumber)
local on = false
function updateBlink()
  local r,g,b = (on and 1 or 0), 0, 0
  for i=0, ledNumber-1 do leds:set(i, r, g, b) end
  on = not on
end
blinkTimer = Timer.new(0.5, updateBlink)  -- 0.5 секунды
blinkTimer:start()
\end{lstlisting}
\textbf{Пример отложенного вызова (через 3 секунды):}
\begin{lstlisting}[style=lua]
local ledNumber = 25
local leds = Ledbar.new(ledNumber)
Timer.callLater(3, function ()
  for i=0, ledNumber-1 do leds:set(i, 0, 0, 1) end  -- включить синий
end)
\end{lstlisting}

\subsection{Событийная модель}
\begin{itemize}
  \item \verb|function callback(event)| — обработчик системных событий.
  \item Примеры: \verb|Ev.COPTER_LANDED| (посадка), \verb|Ev.LOW_VOLTAGE2| (низкое напряжение), \verb|Ev.POINT_REACHED|, \verb|Ev.ALTITUDE_REACHED|.
  \item Внутри \verb|callback| быстро реагируйте на событие; не используйте \verb|sleep()|.
\end{itemize}
\textbf{Примеры:}
\begin{lstlisting}[style=lua,caption={Посадка — выключить индикацию}]
local ledNumber = 25
local leds = Ledbar.new(ledNumber)
function callback(event)
  if (event == Ev.COPTER_LANDED) then
    for i=0, ledNumber-1 do leds:set(i, 0, 0, 0) end
  end
end
\end{lstlisting}
\begin{lstlisting}[style=lua,caption={Низкое напряжение — аварийный красный}]
local ledNumber = 25
local leds = Ledbar.new(ledNumber)
mainTimer = Timer.new(0.1, function () end)
mainTimer:start()
function callback(event)
  if (event == Ev.LOW_VOLTAGE2) then
    mainTimer:stop()
    Timer.callLater(1, function ()
      for i=0, ledNumber-1 do leds:set(i, 1, 0, 0) end
    end)
  end
end
\end{lstlisting}

\section{Анализ подходов к программированию}
\subsection{Конечный автомат}
Логика разделяется на состояния и переключается по событиям (\hyperlink{term:automaton}{конечный автомат}). Это удобно, когда поведение зависит от внешних сигналов: посадка, достижение точки, низкое напряжение. Пример — \verb|4_светодиода.lua|.
\subsection{Блокирующее выполнение}
Использование \verb|sleep()| блокирует систему и мешает другим задачам. Лучше заменять задержки на таймеры. Пример для разбора — \verb|Светодиоды_и_таймер.lua| (как делать не стоит) и корректные таймеры ниже.
\subsection{Асинхронные таймеры}
Эталонный подход — анимации на таймерах без блокировок: вся логика выполняется малыми шагами «на тиках». Пример — \verb|nice_color (1).lua|.

\section{Разбор примеров}
\subsection{4\_светодиода.lua — конечный автомат}
\begin{lstlisting}[style=lua,caption={4_светодиода.lua}]
-- создание порта управления светодиодом
local ledbar = Ledbar.new(25)

-- переменная текущего состояния
local curr_state = "START"

-- таблица функций, вызываемых в зависимости от состояния
action = {
	["START"] = function(x)

		ledbar:set(0, 1, 0.0, 0.0)

		ledbar:set(1, 0.0, 1, 0.0)

		ledbar:set(2, 0.0, 0.0, 1)

		ledbar:set(3, 1, 1, 0.0)

		-- выключение двигателей и конец программы
		ap.push(Ev.ENGINES_DISARM)
		curr_state = "NONE"

	end,
}

-- функция обработки событий, автоматически вызывается автопилотом
function callback(event)
	if (event == Ev.ALTITUDE_REACHED) then
		action[curr_state]()
	end

	if (event == Ev.POINT_REACHED) then
		action[curr_state]()
	end

	if (event == Ev.COPTER_LANDED) then
		sleep(2)
		action[curr_state]()
	end
end

-- вызов функции из таблицы состояний, соответствующей первому состоянию
action[curr_state]()
\end{lstlisting}
\subsection{Светодиоды\_и\_таймер.lua — почему sleep() плохо}
\begin{lstlisting}[style=lua,caption={Светодиоды_и_таймер.lua}]
-- создание порта управления светодиодом
local ledbar = Ledbar.new(8)

-- переменная текущего состояния
local curr_state = "START"

-- таблица функций, вызываемых в зависимости от состояния
action = {
	["START"] = function(x)

		ledbar:set(0, 1, 0.0, 0.0)

		sleep(1000 / 1000.0) 
		ledbar:set(1, 0.0, 1, 0.0)

		sleep(1000 / 1000.0) 
		ledbar:set(2, 0.0, 0.0, 1)

		sleep(1000 / 1000.0) 
		ledbar:set(3, 1, 1, 0.0)

		-- выключение двигателей и конец программы
		ap.push(Ev.ENGINES_DISARM)
		curr_state = "NONE"

	end,
}

-- функция обработки событий, автоматически вызывается автопилотом
function callback(event)
	if (event == Ev.ALTITUDE_REACHED) then
		action[curr_state]()
	end

	if (event == Ev.POINT_REACHED) then
		action[curr_state]()
	end

	if (event == Ev.COPTER_LANDED) then
		sleep(2)
		action[curr_state]()
	end
end

-- вызов функции из таблицы состояний, соответствующей первому состоянию
action[curr_state]()
\end{lstlisting}
\subsection{nice\_color (1).lua — асинхронные таймеры}
\begin{lstlisting}[style=lua,caption={nice_color (1).lua}]
-- https://learnxinyminutes.com/docs/ru-ru/lua-ru/ ссылка для быстрого ознакомления с основами языка LUA

-- Скрипт реализует мигание светодиодов на базовой плате случайными цветами

-- Упрощение вызова функции распаковки таблиц из модуля table
local unpack = table.unpack
-- Количество светодиодов на базовой плате
local ledNumber = 25
-- Создание порта управления светодиодами
local leds = Ledbar.new(ledNumber)

-- Функция, изменяющая цвет 4-х RGB светодиодов на базовой плате
local function changeColor(col)
    for i=0, ledNumber - 1, 1 do
        leds:set(i, unpack(col))
    end
end

-- Функция, реализующая выключение таймера и изменение цвета светодиодов на красный
local function emergency()
    -- Остановка таймера
    timerRandomLED:stop()
    -- Изменение цвета светодиодов (через секунду - т.к. после остановки таймера timerRandomLED
    -- его функция выполнится еще раз)
    Timer.callLater(1, function () changeColor({1, 0, 0}) end)
end

-- Функция обработки событий, автоматически вызывается автопилотом
function callback(event)
    -- Вызов функции emergency() при низком напряжении на аккумуляторе
    if (event == Ev.LOW_VOLTAGE2) then
        emergency()
    end
end

color = {0, 0, 0} 
step = 0.05 
index = 1

-- Создание таймера, каждую секунду меняющего цвета каждого из 4-х светодиодов на случайные
timerRandomLED = Timer.new(0.05, function ()
    color[index] = color[index] + step
    if color[index] > 0.95 or color[index] < 0.05 then
        if index < 3 then
            index = index + 1
        else
            index = 1
            step = -step
        end
    end
    changeColor(color)
end)
-- Запуск созданного таймера
timerRandomLED:start()
\end{lstlisting}

\section{Задания для лабораторной работы}
Используйте асинхронные таймеры, глобальные переменные для таймеров и состояний, полный проход по индексам \verb|0..24|. Формулировки ниже предельно простые; рядом указаны примеры для ориентира.
\subsection{Базовые алгоритмы}
\begin{enumerate}
  \item \textbf{Зелёный свет.} Включите \verb|{0,1,0}| на всех индексах \verb|0..24|. Подсказка: пример инициализации в разделе «Адресация и буфер».
  \item \textbf{Отложенный запуск.} Через 3 секунды включите \verb|{0,0,1}|, до этого — \verb|{0,0,0}|. Подсказка: используйте \verb|Timer.callLater|.
  \item \textbf{Мигание.} Каждые 0.5с переключайте \verb|{1,0,0}| и \verb|{0,0,0}|. Подсказка: глобальный таймер и булево состояние; см. «Асинхронные таймеры».
  \item \textbf{Бегущий огонь.} На каждом тике добавляйте один зелёный индекс, остальные — выключены. Подсказка: счётчик позиции.
  \item \textbf{Светофор.} Красный → жёлтый → зелёный по кругу. Подсказка: конечный автомат; см. «4\_светодиода.lua».
\end{enumerate}
\subsection{Алгоритмические задачи}
\begin{enumerate}[resume]
  \item \textbf{Пульсация (PWM).} Плавно меняйте яркость красного от \verb|0.0| до \verb|1.0| и обратно. Подсказка: переменные \verb|value| и \verb|step|.
  \item \textbf{Стробоскоп.} Быстро мигающий белый цвет. Подсказка: период таймера \verb|0.03|с.
  \item \textbf{Чётность индекса.} Чётные — красный, нечётные — синий для всех \verb|0..24|. Подсказка: \verb|i % 2|.
  \item \textbf{Полицейская мигалка.} Меняйте местами «синий/красный» на чётных/нечётных каждые 0.5с. Подсказка: два шаблона и переключение.
  \item \textbf{Генератор шума.} Случайный цвет для каждого индекса. Подсказка: \verb|math.random()| и полный проход буфера.
\end{enumerate}
\subsection{Работа с событиями}
\begin{enumerate}[resume]
  \item \textbf{Посадка.} При \verb|Ev.COPTER_LANDED| погасите все индексы. Подсказка: обработчик \verb|callback|.
  \item \textbf{Низкое напряжение.} При \verb|Ev.LOW_VOLTAGE2| остановите таймеры и включите красный. Подсказка: \verb|timer:stop()| и \verb|Timer.callLater|.
  \item \textbf{Loading bar.} Заполняйте индексы от \verb|0| до \verb|24|, удерживая включённые. Подсказка: счётчик \verb|k|.
  \item \textbf{SOS.} Белые импульсы по схеме «••• ——— •••». Подсказка: таблица длительностей и автомат.
  \item \textbf{Бинарный счётчик.} Отобразите биты секунд (\verb|os.time()%32|) на линейке. Подсказка: проверка каждого бита.
\end{enumerate}

\end{document}

